\documentclass[aspectratio=169]{beamer}
\usetheme{Madrid}
\usepackage[utf8]{inputenc}
\usepackage[T1]{fontenc}
\usepackage{tikz}
\usepackage{amssymb}
\usepackage{graphicx}
\setbeamertemplate{navigation symbols}{}
\setbeamercolor{title}{fg=white,bg=blue!65!cyan}
\setbeamercolor{frametitle}{fg=white,bg=blue!65!cyan}
\setbeamercolor{block title}{bg=blue!65!cyan!20,fg=black}
\setbeamercolor{block body}{bg=blue!65!cyan!8,fg=black}
\setbeamercolor{item}{fg=blue!65!cyan}
\setbeamercolor{itemize item}{fg=blue!65!cyan}
\setbeamertemplate{itemize item}[triangle]

\title{Lists}
\subtitle{Storing many things under one name — index and .append() only}
\author{ByteForge --- Grade 8 --- App Builder}
\date{}

\begin{document}

\begin{frame}
    \titlepage
\end{frame}

\begin{frame}{Today's Plan}
    \begin{itemize}
        \item Your First List
        \item[$\bigstar$] BOSS --- Three Items, One List
    \end{itemize}
\end{frame}

\begin{frame}{A list holds MANY things under one name}
    \large
    \begin{itemize}
        \item A direct comparison to Grade 7's arrays
        \item \texttt{inventory = []} starts empty, ready to hold as many items as needed.
    \end{itemize}
\end{frame}


\begin{frame}{\texttt{.append()} adds to the END}
    \large
    \begin{itemize}
        \item \texttt{inventory.append("sword")} genuinely adds a new item after the list already exists
        \item The list grows one call at a time.
    \end{itemize}
\end{frame}


\begin{frame}{Indexing starts at 0, not 1}
    \large
    \begin{itemize}
        \item \texttt{inventory[0]} reads the FIRST item ever added; the THIRD item added sits at index 2, not 3
        \item An off-by-one trap just like \texttt{range()}'s.
    \end{itemize}
\end{frame}

\begin{frame}{Real-World Examples}
    \large
    \begin{itemize}
        \item Every shopping cart, playlist, or to-do list app stores its items in exactly this shape
        \item An ordered list that grows one \texttt{.append()}-equivalent call at a time, read back by position.
    \end{itemize}
\end{frame}


\begin{frame}{Let Them Discover}
    \Large
    Add three real items to a list on paper (write them in order). Which index number is the SECOND item you added? Is it 1 or 2 — and why?
\end{frame}

\begin{frame}{Your Turn}
    \Large
    Open the student map, find \textbf{Lists}, and work through it
    hands-on --- everything you just saw, they now do themselves.
\end{frame}

\end{document}
